\subsection{Source Resolution With Independent Phasor Errors}
\label{app:phasor-noise-contrast}

We tested whether the loss on noisy records begins at the source
decision or at the alarm. The observation $y_k$ contains 31 terminal
coordinates from the eight PMUs. These comprise 16 log magnitudes
of voltage and current, eight unwrapped voltage angles and seven
usable current angles in radians. Given a possible onset $u$, let $b_0,b_1,b_2$
denote the means of three consecutive $h$-report blocks. The first two
end before $u$ and the last starts at $u$. We use
\[
 d_h=b_2-2b_1+b_0,\qquad
 C_h(q)=C_{h,0}+(6q/h)I .
\]
The contrast cancels a constant or linear time trend. Independent
errors of variance $q$ per polar coordinate add $6q/h$ to its
variance because the three disjoint block means have squared weights
$1,4,1$. Eight quiet calibration records from the reference plant
supply the background mean and covariance $C_{h,0}$ for each $h$.

To estimate $q$ from a new record, we form 100 nonoverlapping
three-report second differences from its first 300 reports and divide
each by $\sqrt{6}$. Let $\lambda_j$ be the eigenvalues of the
calibration covariance of these differences. Let $m_j$ be the
startup mean square in eigenvector $j$ after subtracting the
calibration mean. We fit $q\geq0$ by minimizing the Gaussian score
\[
 \sum_{j=1}^{31}\left[\log(\lambda_j+q)+
                   \frac{m_j}{\lambda_j+q}\right].
\]
This is a marginal score for the differences. Their background can
remain correlated in time. For source $a$ and load severity $s$, the
nonlinear second plant provides an isolated response $\mu_{a,s}$ in
contrast coordinates. We evaluate the Gaussian likelihood with mean
$\mu_{a,s}$ plus the background mean and covariance $C_h(q)$.
The source decision maximizes posterior mass after summing over four
equally likely load sizes ($-10,-3,3,10\%$). Normal has prior mass
$1/2$. Load and line each have mass $1/4$, divided equally among
their nine candidate assets. This decision minimizes source-label
error under the specified mixture for a given onset and fitted
covariance.

\begin{table}[t]\centering\small
\caption{Source decisions one second after a supplied onset. Each
condition uses the same 432 load, 108 line and 540 quiet windows.
Quiet$^+$ counts quiet windows labeled as physical events.}
\label{tab:pmu-contrast-window}
\begin{tabular}{lrrr}\hline
Covariance & Load /432 & Line /108 & Quiet$^+$ /540\\\hline
\multicolumn{4}{l}{0.05\% added RMS TVE}\\
Nominal & 256 & 108 & 536\\
Diagonal, estimated & 413 & 108 & 14\\
Full, estimated & 409 & 108 & 0\\
Full, supplied & 409 & 108 & 0\\
\multicolumn{4}{l}{0.10\% added RMS TVE}\\
Nominal & 182 & 108 & 538\\
Diagonal, estimated & 406 & 108 & 4\\
Full, estimated & 404 & 108 & 0\\
Full, supplied & 404 & 108 & 0\\
\hline\end{tabular}

\end{table}

For Table~\ref{tab:pmu-contrast-window}, each of 12 separate quiet
test records from the second plant supplies 90 disjoint windows.
In 45 per record we add an isolated simulated response in the same
coordinates as the classifier's templates. The other 45 contain
background alone. Fresh independent Gaussian errors in log magnitude
and phase have standard deviation equal to the nominal added RMS
total vector error divided by $\sqrt{2}$ (small-error approximation).
We froze the classifier before generating the windows. The primary
readout uses $h=30$ at 30 reports/s. All four arms see the same
windows. The diagonal arm discards between-channel covariance after
estimating $q$. The supplied arm receives the injected $q$.
At 0.05\%, supplying $q$ gave the same counts as estimating it.
The diagonal fit recovered four more loads but labeled 14 quiet
windows as events. Compared with fixed nominal covariance, the
estimated full covariance recovered 10--16 additional loads in each
of the 12 background records. Windows within a record share its
startup estimate.

\begin{table}[t]\centering\small
\caption{Continuous replay at 0.05\% added RMS total vector error.
Det., Fam.\ and Exact count detection, correct family and correct
family and asset before restoration among 72 interventions per plant.
Extra counts unmatched event-record alarms. Quiet$^+$ counts alarms
in \ContrastQuietHours{} normal hours per row.}
\label{tab:pmu-contrast-stream}
\begin{tabular}{llrrrrr}\hline
Plant & Covariance & Det. & Fam. & Exact & Extra & Quiet$^+$\\\hline
F01 & Nominal & 0 & 0 & 0 & 12 & 12\\
F01 & Estimated & 12 & 10 & 3 & 36 & 2\\
P2 & Nominal & 0 & 0 & 0 & 12 & 12\\
P2 & Estimated & 16 & 13 & 6 & 84 & 2\\
\hline\end{tabular}

\end{table}

Table~\ref{tab:pmu-contrast-stream} uses all 24 archived noisy records
per plant, with the same 12 six-intervention schedules and 12 quiet
records across plants. Both replays used the frozen isolated responses
from the second plant, including the replay on the reference plant.
The unknown-onset rule slides the three blocks
one report at a time after 12~s of startup. It publishes after the
same source exceeds the source-to-normal odds threshold for six
consecutive reports and rearms after six consecutive reports below
threshold. The common log threshold was the greater of $\log 99$ and
the highest estimated-covariance score on the eight quiet calibration
records with or without fresh 0.05\% noise (plus a numerical margin
on the calibration maximum). The first unmatched alarm during each
intervention determines its diagnosis. All remaining alarms count as
extras. Sharing a threshold permits the two arms to incur different
false-alarm rates. In this archived stress test each complex phasor is
multiplied by $1+(0.0005/\sqrt{2})(e_1+\mathrm{i}e_2)$, with fresh
independent standard normal components at each report and terminal.

On these same noisy trajectories we also read the unconstrained source
decision 30 reports from each true onset. With estimated covariance it
identified \ContrastFAtOnset{}/72 sources on the reference plant and
\ContrastPAtOnset{}/72 on the second plant. On each plant it identified
10 of 12 first interventions and 22 of 48 that followed restoration
by less than 3~s. Requiring the classifier
to find the onset reduced those counts to \ContrastFExact{} and
\ContrastPExact{}. On the second plant the clock also produced
\ContrastPExtra{} unmatched event-record alarms and \ContrastPQuiet{}
quiet alarms. These trajectories include continuously sampled load
sizes and the residual motion of previous interventions, whereas the
isolated-window experiment uses the same responses for insertion and
inference. The comparison does not isolate the contribution of each
difference. The main history bank carries the estimated state through
preceding interventions.

We tested the discrete load prior on 36 newly simulated two-event
sequences from the second plant. Each load or line family order
occurred nine times. Load size was uniform in absolute value from
3\% to 14\%, with either sign equally likely. The gap after restoration
was 0.5--3~s. Each sequence had its own instrument noise and independent
0.05\% RMS phasor error. We froze both rules before generating the
observations. Both saw the same eight PMUs and used the same contrast,
covariance and alarm threshold. For the continuous prior we drew a
secant through the 3\% and 10\% responses separately for each sign,
extended it to 14\%, and integrated the Gaussian likelihood along it.
For endpoints $l,r$ in whitened coordinates and observation $z$,
the likelihood relative to the null, conditional on the sign, is
\[
 \int_0^1 \exp\!\left(z^\top m(u)-\tfrac12\|m(u)\|^2\right)du,
 \qquad m(u)=l+u(r-l).
\]
Completing the square reduces this integral to the difference of two
normal distribution functions. We average over sign before deciding
the asset.

With the onset supplied, integration increased exact first-event
diagnoses from \SizeGridFirst{} to \SizeIntegralFirst{} out of 36.
For the second intervention it increased them from \SizeGridSecond{}
to \SizeIntegralSecond{} out of 36. With the alarm clock running,
both rules detected every first intervention and missed every
second one before restoration. Exact diagnoses fell from
\SizeGridExact{} to \SizeIntegralExact{} out of 72, and unmatched
event-record alarms increased from \SizeGridExtra{} to
\SizeIntegralExtra{}. Neither rule raised a quiet alarm in 1.44~h
from 18 additional simulated records. A more accurate size model
helped when the first switch time was given. It did not supply the
recovery prediction needed to interpret successive switches.
